%\newcommand{\refTeorem}[1]{(teorem \ref{#1})}
%\newtheorem{theorem}{Teorem}
%\begin{theorem}[Newtons metode]
%\label{newtonsmetode}
%\[x_{n+1} = x_n - \frac{f(x_n)}{f'(x_n)}\]
%\end{theorem}


\oppgave{1}
\oppgaveDelStart
\oppgaveDel{a}

Newtons metode:
\begin{equation}
\label{newton}
\begin{split}
&\text{1. Gjett at $x_0$ er et nullpunkt.}\\
&\text{2. Gjør følgende steg om og om igjen:}\\
&x_{n+1} = x_n - \frac{f(x_n)}{f'(x_n)}
\end{split}
\end{equation}

\vspace{0.5cm}

Tilnærmet verdi for $\sqrt[3]{2}$ ved bruk av Newtons metode \eqref{newton}:
\begin{align*}
&x = \sqrt[3]{2}&\\
&x^3 = 2\\
&x^3-2 = 0\\
&f(x) = x^3 - 2, \ \ f'(x) = 3x^2\\\\
%&\boxed{x_{n+1} = x_n - \frac{x^3_n - 2}{3x^2_n}}\\\\
&x_1 = x_0 - \frac{f(x_0)}{f'(x_0)}\\
&\text{Vi velger } x_0 = 1\\
&x_1 = 1 - \frac{1^3 - 2}{3(1)^2}\\
&x_1 = 1.333333\dots\\\\
&x_2 = 1.33333 - \frac{1.33333^3 - 2}{3(1.33333)^2}\\
&x_2 = 1.263889\\\\
&x_3 = 1.263889 - \frac{1.263889^3 - 2}{3(1.263889)^2}\\
&x_3 = \uuline{1.26003}
\end{align*}

%\vspace{1cm}
\newpage
Tilnærmet verdi for $\sqrt[3]{2}$ ved bruk av Newtons metode i Python:
\vspace{0.5cm}
\begin{lstlisting}[language=Python]
import math;

val = 2**(1/3)

def f(x):
    return x**3 - 2

def Df(x):
    return 3*x**2

feil =  0.0001
x = 1
it = 0

while (abs(x - val) > feil):
    print(f'{it} x = {round(x,8)}, f(x) = {round(f(x),8)}')
    x = x - f(x)/Df(x)
    it += 1
print(f'2**(1/3) = {val}')

\end{lstlisting}

\newpage

\oppgaveDel{b}

Lineær tilnærming:
\begin{equation}
\label{linear}
\begin{split}
&\text{Hvis $f(x)$ er deriverbar i $x=a$, da er}\\[0.1cm]
&L(x) = f'(a)(x-a) + f(a)\\[0.1cm]
&\text{en lineær tilnærming til $f(x)$ i $x=a$.}
\end{split}
\end{equation}

\vspace{0.5cm}

Tilnærmet verdi for $\sqrt[3]{7,96}$ ved bruk av lineær tilnærming \eqref{linear}:
\begin{align*}
&\text{La } f(x) = \sqrt[3]{x} = x^{\frac{1}{3}}\\
&\text{Da er } f'(x) = \frac{1}{3} x^{- \frac{2}{3}} = \frac{1}{3 \sqrt[3]{x^2}}\\\\
&\text{Det nærmeste tallet til $7.96$ som har en perfekt kubikkrot er $8$ ($\sqrt[3]{8} = 2$).}\\
&\text{Vi lar } x = 8\\
&f(8) = \sqrt[3]{8} = 2\\
&f'(8) = \frac{1}{3 \sqrt[3]{8^2}} = \frac{1}{3 \sqrt[3]{64}} = \frac{1}{3(4)} = \frac{1}{12}\\\\
&\text{Dette setter vi inn i formelen for lineær tilnærming \eqref{linear}:}\\
&a = 8, \ x = 7,96, \ x - a = 7.96 - 8 = -0.04\\
&\sqrt[3]{7,96} \approx 2 + \frac{1}{12} \cdot (-0.04)\\
&\sqrt[3]{7,96} \approx \uuline{1.996667}\\\\
&\text{Er dette overestimat eller underestimat?}\\
%&\text{Først finner vi $f''(x)$:}\\
&f''(x) = \frac{1}{3} (- \frac{2}{3} x^{-\frac{5}{3}}) = - \frac{2}{9} x^{-\frac{5}{3}} \ \ \Rightarrow \ \ f''(x) < 0\\
&\text{Vi ser at $f''(x)$ er negativ. Funksjonen er konkav.}\\
&\text{Dette tilsier at at det er et overestimat.}
\end{align*}

\newpage

\oppgaveDel{c}

Teorem for derivasjon av inverse funksjoner:
\begin{equation}
\label{invers}
&(h^{-1})'(y) = \frac{1}{h'(h^{-1}(y))}
\end{equation}

\vspace{0.5cm}

Finn $(h^{-1})'(13)$ for &h(x) = x^3 + 5$ uten å bestemme $h^{-1}(x)$:
\begin{align*}
&h(x) = 13\\
&x^3 + 5 = 13\\
&x^3 = 8\\
&x = 2\\
\Rightarrow \ &h^{-1}(13) = 2\\\\
%&\text{Videre, finner vi } h'(x):\\
&h'(x) =  3x^2\\
&h'(h^{-1}(13)) = h'(2) = 3(2)^2 = 12\\\\
&\text{Vi setter $h'(h^{-1}(13))$ inn i \eqref{invers}:}\\
&(h^{-1})'(13) = \frac{1}{h'(h^{-1}(13))} = \frac{1}{h'(2)} = \frac{1}{12}
\end{align*}

\oppgaveDelSlutt

% % % % % % % % % % % % % % % % % % % % % % % % % % % % % % % % % % % % % % % %
\newpage
\oppgave{2}
\oppgaveDelStart

\oppgaveDel{a}
\begin{align*}
&\int_{0}^{1} \Big(2 + 6x^2 + 3e^x\Big) \ dx = \Big[2x + 2x^3 + 3e^x\Big]_0^1 = (2+2+3e) - (0+0+3) = \uuline{1+3e}
\end{align*}

\oppgaveDel{b}
\begin{align*}
&\int_1^3 \Big(\frac{2x^2 + 2x + 3}{x}\Big) \ dx = \int_1^3 \Big(2x+2+\nicefrac{3}{x}\Big) \ dx = \Big[ x^2 + 2x + 3 \ln |x| \Big]_1^3 = \uuline{12 + 3 \ln 3}
\end{align*}

\oppgaveDel{c}
\begin{align*}
%&\int x e^{7x^2} \ dx\\
&u = 7x^2 \ \ \Rightarrow \ \ \frac{du}{dx} = 14x \ \ \Rightarrow \ \ du = 14x \ dx \ \ \Rightarrow \ \ x \ dx = \frac{du}{14}\\\\
&\int \Big( x e^{7x^2} \Big) \ dx = \frac{1}{14} \int e^u \ du = \frac{1}{14} \ e^u + C = \uuline{\frac{1}{14} \ e^{7x^2} + C}\\
\end{align*}

\oppgaveDel{d}
\begin{align*}
%&I = \int_{\frac{1}{e}}^e \frac{1}{x((\ln x)^2 + 1)} \ dx\\\\
& u = \ln x \ \ \Rightarrow \ \ \frac{du}{dx} = \frac{1}{x} \ \ \Rightarrow \ \ du = \frac{dx}{x} = \frac{1}{x} dx\\
&x_1 = \frac{1}{e} \ \ \Rightarrow \ \ u_1 = \ln(\nicefrac{1}{e}) = \ln 1 - \ln e = 0-1 = -1\\
&x_2 = e \ \ \Rightarrow \ \ u_2 = \ln e = 1\\\\
&\int_{\nicefrac{1}{e}}^e \bigg( \frac{1}{x \big( (\ln x)^2 + 1 \big) } \bigg) \ dx = \int_{u_1 = -1}^{u_2 = 1} \bigg( \frac{1}{u^2 + 1} \bigg) \ du\\\\
%&\int \frac{1}{u^2 + 1} \ du = \tan^{-1} u + C\\\\
&\int_{-1}^1 \bigg( \frac{1}{u^2+1} \bigg) \ du = \Big[ \tan^{-1} u \Big]_{-1}^{1} = \tan^{-1} (1) - \tan^{-1} (-1) = \frac{\pi}{4} - (- \frac{\pi}{4}) = \uuline{\frac{\pi}{2}}
\end{align*}

\newpage


Delvis integrasjon:
\begin{equation}
\label{delvisIntegrasjon}
&\int u \ dv = uv - \int v \ du
\end{equation}
%\vspace{0.5cm}
\oppgaveDel{e}
Vi bruker delvis integrasjon \eqref{delvisIntegrasjon} for å løse integralet:
\begin{align*}
%&\int x e^{-2x} \ dx\\
&u = x, \ \ dv = e^{-2x} \ dx\\
& du = dx, \ \ \int dv = \int e^{-2x} \ dx\\
%&\text{For integralet } \int e^{-2x} \ dx\\
%&\text{La } u = -2x \Rightarrow du = -2 \ dx \Rightarrow - \frac{1}{2} \ du = dx\\
&\int e^{-2x} = \int e^u - \frac{1}{2} \ du = - \frac{1}{2} \int e^u \ du = - \frac{1}{2} e^u + C = -\frac{1}{2} e^{-2x} + C\\
\Rightarrow \ \ &v = - \frac{1}{2} e^{-2x}
\end{align*}
\begin{align*}
\int x e^{-2x} \ dx &= x(-\frac{1}{2} e^{-2x}) - \int (- \frac{1}{2} e^{-2x}) \ dx\\
&= - \frac{1}{2} x e^{-2x} + \frac{1}{2} \int e^{-2x} \ dx\\
&= - \frac{1}{2} x e^{-2x} + \frac{1}{2} (- \frac{1}{2} e^{-2x}) + C\\
&= - \frac{1}{2} x e^{-2x} - \frac{1}{4} e^{-2x} + C
\end{align*}

\oppgaveDelSlutt

% % % % % % % % % % % % % % % % % % % % % % % % % % % % % % % % % % % % % % % %

\newpage
\oppgave{3}
%$$f(x) = e^{-\frac{x}{2}}, \quad g(x) = 2e^x$$
\oppgaveDelStart
\oppgaveDel{a}
\begin{align*}
&u = - \frac{x}{2} = - \frac{1}{2} x \ \ \Rightarrow \ \ du = - \frac{1}{2} \ dx \ \ \Rightarrow \ \ dx = -2 \ du
\end{align*}
\begin{align*}
\int_0^{\ln 2} f(x) \ dx &= \int_0^{\ln 2} e^{- \frac{x}{2}} \ dx = \int_0^{\ln 2} e^u (-2 \ du) = -2 \int_0^{\ln 2} e^u \ du\\\\
&= -2 \Big[ e^u \Big]_0^{\ln 2} = -2 \Big[ e^{-\frac{x}{2}} \Big]_0^{\ln 2} = -2e^{-\frac{\ln2}{2}} - (-2e^0)\\\\
&= \uuline{2 - 2e^{\ln (2^{-\nicefrac{1}{2}})}}
\end{align*}

\oppgaveDel{b}
\begin{align*}
&e^{-\frac{x}{2}} = 2e^x\\
&\frac{e^{-\frac{x}{2}}}{e^x} = 2\\
&e^{-\frac{x}{2}-x} = 2\\
&e^{-\frac{3}{2}x} = 2\\
&\ln e^{-\frac{3}{2}x} = \ln 2\\
&-\frac{3}{2}x \ln e = \ln 2\\
&-\frac{3}{2}x = \ln 2\\
&x = - \frac{2}{3} \ln 2 = - \frac{2}{3}(0.6931) = -0.4621\\
\Rightarrow \ & \text{$f$ og $g$ møtes ved $x = a = -0.4621$}\\\\
%&\text{$y$-akse ($x=0$)}\\
&f(0) = e^0 = 1\\
&g(0) = 2e^0 = 2\\
&\text{Ved $x=0$ er $g(x) > f(x)$}\\
\Rightarrow \ &\text{Kurven $g$ er over $f$}
\end{align*}
\begin{align*}
%&\int 2e^x \ dx = 2 e^x\\
%&\int e^{-\frac{x}{2}} \ dx = -2 e ^{-\frac{x}{2}}\\
A &= \int_{x=a}^0 \left[ g(x) - f(x) \right] \ dx = \int (2e^x - e^{-\frac{x}{2}}) \ dx\\
&= \left[ 2e^x - (-2 e ^{-\frac{x}{2}}) \right]_a^0 = (2e^0 + 2e^0) - (2e^{0.4621} + 2e^{\frac{0.4621}{2}})\\
&= 4 - 3.77976\\
A &= \uuline{0.22024}
\end{align*}
\oppgaveDelSlutt

% % % % % % % % % % % % % % % % % % % % % % % % % % % % % % % % % % % % % % % %

\newpage
\oppgave{4}
\oppgaveDelStart
\oppgaveDel{a}
\begin{align*}
&f(3) = 4 + \int_3^3 \sin \left[ \ln(t^3+2t-\frac{1}{t}) \right] \ dt = 4 + 0 = \uuline{4}
\end{align*}
\oppgaveDel{b}
\begin{align*}
&\text{Vi bruker fundamentalteoremet i kalkulus:}\\
&F(x) = \int_a^x f(t) \ dt \Rightarrow F'(x) = f(x)\\
&f'(x) = \sin \left[ \ln(x^3+2x-\frac{1}{x}) \right]
\end{align*}
\oppgaveDelSlutt

% % % % % % % % % % % % % % % % % % % % % % % % % % % % % % % % % % % % % % % %

\oppgave{5}
\oppgaveDelStart
\oppgaveDel{a}
\begin{align*}
&\lim_{x\to\infty} xe^{-x} = \lim_{x\to\infty} \frac{x}{e^x} = \uuline{0}
\end{align*}

\oppgaveDel{b}
\begin{align*}
%&\lim_{x \to 0} \sin^x x\\\\
&y = (\sin x)^x \ \ \Rightarrow \ \ \ln y = x \ln (\sin x)\\\\
&\text{Først finner vi grenseverdien:}\\
&L = \lim_{x\to0} x \ln(\sin x) = x \ln(\sin x) = \frac{\ln(\sin x)}{\frac{1}{x}}\\
&x\to 0^+, \ L = \frac{-\infty}{+\infty}\\\\
&\text{Derfor bruker vi L'Hôpitals regel.}\\
&\frac{d}{dx} \left[ \ln(\sin x) \right] = \frac{\cos x}{\sin x} = \cot x\\
&\frac{d}{dx} \left[ \frac{1}{x} \right] = - \frac{1}{x^2}\\
&L = \lim_{x\to0} x \ln(\sin x) = \lim_{x\to0} \frac{\ln(\sin x)}{\frac{1}{x}} = \lim_{x\to0} \frac{\cot x}{- \frac{1}{x^2}} = \lim_{x\to 0} -x^2 \ \cot x = 0\\\\
&\Rightarrow \ln y = 0 \Rightarrow y = e^0 = 1\\
&\lim_{x\to 0} y = \lim_{x \to 0} 1 = 1 \Rightarrow \lim_{x\to 0} \sin^x x = \uuline{1}
\end{align*}
\oppgaveDelSlutt

% % % % % % % % % % % % % % % % % % % % % % % % % % % % % % % % % % % % % % % %

\newpage
\oppgave{6}
\oppgaveDelStart
\oppgaveDel{a}
\begin{align*}
&\text{Vi bruker pytagorassetningen:}\\
&c^2 = a^2 + b^2\\
&x^2 + y^2 = 5^2 = 25\\
&\frac{dx}{dt} = 0,4\ m/s\\
&\text{Vi må finne $\frac{dy}{dt}$ når $x = 3m$}\\
&x = 3 \Rightarrow 3^2 + y^2 = 25\\
&y^2 = 25-9 = 16\\
&y = 4\\\\
&2x \frac{dx}{dt} + 2y \frac{dy}{dt} = 0\\
&\Rightarrow x \frac{dx}{dt} + y \frac{dy}{dt} = 0\\
&\text{Men }\frac{dx}{dt} = 0,4\ m/s, \ x = 3, \ y =4\\
&\text{så } x \frac{dx}{dt} + y \frac{dy}{dt} = 0 \Rightarrow 3(0,4) + 4 \frac{dy}{dt} = 0\\
&\Rightarrow 1,2 + 4 \frac{dy}{dt} = 0\\
&\Rightarrow \frac{dy}{dt} = \uuline{-0.3 m/s}
\end{align*}

\oppgaveDel{b}
\begin{align*}
&A = \frac{1}{2} x y\\
&\frac{dA}{dt} = \frac{1}{2} \left( x \frac{dy}{dt} + y \frac{dx}{dt} \right)\\
&x = 3, \ \frac{dy}{dt} = -0,3m/s\\
&y = 4, \ \frac{dx}{dt} = 0.4 m/s\\
&\frac{dA}{dt} = \frac{1}{2} \left( 3 (-0,3) + 4 (0,4) \right) = \frac{1}{2}(0,7)\\\\
&\frac{dA}{dt} = \uuline{0,35 \ m/s}
\end{align*}
\oppgaveDelSlutt

% % % % % % % % % % % % % % % % % % % % % % % % % % % % % % % % % % % % % % % %
